% ============================================================================
% Fichier  : partie5/P05_C24_investigation.tex
% Titre    : Premiers Pas en Investigation Numérique
% Auteur   : MINKA MI NGUIDJOÏ Thierry Emmanuel
% Version  : 1.0 -- novembre 2026
% ============================================================================
% Base : brouillon « Livre VIII » (investigation numérique, chaîne de
% garde). Sources : ISO/IEC 27037:2012 ; RFC 3227 (2002) ; NIST SP 800-86
% (2006) ; loi n° 2010/012, art. 52 (3) et 53 (1), texte officiel ; Trilemme
% CRO (Minka Mi Nguidjoi et al., 2025). Renvoi à l'UE SEC4031.
% ============================================================================

\chapter{Premiers Pas en Investigation Numérique}
\label{chap:investigation}
\chaptermark{Investigation numérique}

\epigraph{%
    « Tout contact laisse une trace. »%
}{Edmond Locard}

\niveauChapitre{Intermédiaire}{Chapitres~\ref{chap:crypto-sym},
\ref{chap:droit} et~\ref{chap:operations}.}

\begin{boxobjectifs}
\begin{itemize}
    \item Expliquer pourquoi la réponse aux incidents doit préserver les
          preuves.
    \item Appliquer l'ordre de volatilité.
    \item Décrire les quatre étapes de l'ISO/IEC~27037.
    \item Tenir une chaîne de conservation.
    \item Situer le rôle de l'entreprise, des enquêteurs et de l'expert.
\end{itemize}
\end{boxobjectifs}

\bigskip

% ============================================================================
\section*{Question directrice : que reste-t-il pour prouver ce qui s'est passé~?}
% ============================================================================

LiZbiZ porte plainte. Les enquêteurs demandent les preuves~: la mémoire
des serveurs a été effacée par le redémarrage de 22~h~05, plusieurs
journaux locaux ont été écrasés par la réinstallation des machines, et
personne n'a noté qui a manipulé quoi. Ce chapitre donne les premiers
gestes qui préservent la preuve~; l'investigation proprement dite fait
l'objet de l'UE SEC4031.

\medskip
\textbf{Plan du chapitre.}
\begin{enumerate}
    \item Réponse et preuve (section~\ref{sec:c24-reponse}).
    \item L'ordre de volatilité (section~\ref{sec:c24-volatilite}).
    \item Les quatre étapes de l'ISO/IEC~27037 (section~\ref{sec:c24-27037}).
    \item La chaîne de conservation (section~\ref{sec:c24-chaine}).
    \item L'entreprise, l'enquêteur, l'expert (section~\ref{sec:c24-roles}).
\end{enumerate}

\separateur

% ============================================================================
\section{Réponse et Preuve}
\label{sec:c24-reponse}
% ============================================================================

L'objectif de la réponse est de rétablir~; celui de l'investigation est
d'établir les faits. Les deux se gênent~: réinstaller détruit les traces,
attendre prolonge l'arrêt. La solution est de préserver \emph{avant}
d'éradiquer~: images des disques, capture de la mémoire, copie des
journaux, avec leurs empreintes. Le NIST décrit cette intégration de la
forensique dans la réponse aux incidents~\cite{nist80086}.

% ============================================================================
\section{L'Ordre de Volatilité}
\label{sec:c24-volatilite}
% ============================================================================

La RFC~3227 recommande de collecter d'abord les données qui disparaissent
le plus vite~\cite{rfc3227}.

\begin{tableaucours}{cL{5.6cm}L{5.4cm}}{%
    \tableauHeader{Ordre} &
    \tableauHeader{Donnée} &
    \tableauHeader{Disparaît}
}
1 & Registres, cache & En quelques nanosecondes \\
2 & Mémoire vive, processus, connexions réseau & À l'extinction \\
3 & Fichiers temporaires & Au redémarrage ou au nettoyage \\
4 & Disques & À la réécriture \\
5 & Journaux distants, données de surveillance & Selon leur rétention \\
6 & Configuration physique, topologie & Rarement \\
7 & Supports d'archive & Très rarement \\
\end{tableaucours}
\vspace{-0.8em}
\begin{center}{\small\itshape Tableau~24.1. L'ordre de volatilité (d'après la RFC~3227)}\end{center}

\begin{boxscenario}{Opération GAROUA. Ce que le redémarrage a effacé}
\textbf{Constatation.} À 22~h~05, M.~ABENA redémarre les serveurs de
Garoua et de Douala.

\medskip
\textbf{Analyse} (confiance élevée). Le redémarrage a effacé la mémoire
vive~: les processus de l'attaquant, ses connexions en cours, et
peut-être des clés de chiffrement qu'un examen de la mémoire aurait pu
retrouver. Isoler du réseau sans éteindre aurait préservé ces traces.
\end{boxscenario}

% ============================================================================
\section{Les Quatre Étapes de l'ISO/IEC 27037}
\label{sec:c24-27037}
% ============================================================================

\begin{tableaucours}{L{2.8cm}L{9.4cm}}{%
    \tableauHeader{Étape} &
    \tableauHeader{Principe}
}
Identification & Repérer les supports et les données qui peuvent
constituer une preuve \\
Collecte & Prendre les supports physiques, quand c'est possible et
légitime \\
Acquisition & Produire une copie fidèle, vérifiable par empreinte \\
Préservation & Protéger contre toute altération et tracer chaque accès \\
\end{tableaucours}
\vspace{-0.8em}
\begin{center}{\small\itshape Tableau~24.2. Les étapes de l'ISO/IEC~27037~\cite{iso27037}}\end{center}

\begin{boxprocedure}
\textbf{Acquérir un disque}
\begin{enumerate}
    \item Photographier et décrire le support (marque, numéro de série).
    \item Brancher le disque derrière un bloqueur d'écriture.
    \item Calculer l'empreinte SHA-256 du disque source.
    \item Produire l'image (format brut ou E01).
    \item Calculer l'empreinte de l'image et vérifier qu'elle est identique.
    \item Consigner chaque étape~: date, heure, personne, outil et version.
\end{enumerate}
\end{boxprocedure}

% ============================================================================
\section{La Chaîne de Conservation}
\label{sec:c24-chaine}
% ============================================================================

\begin{boxdef}
\termecle{Chaîne de conservation} (\textit{chain of custody})~: suite
ininterrompue d'enregistrements qui retrace, pour chaque pièce, qui l'a
détenue, quand, où et pour quoi faire, depuis sa collecte jusqu'à sa
production en justice.
\end{boxdef}

Une rupture dans la chaîne (un disque manipulé sans trace, une image
modifiée sans empreinte) permet à la défense de contester la pièce. La
discipline est la même que pour l'auditeur qui remet ses pièces à la
justice (cours INF3226, Chapitre~29).

\begin{boxinfo}
\textbf{Le Trilemme CRO.} Une preuve numérique doit être confidentielle
(elle contient des données sensibles), fiable (elle n'a pas été
altérée) et opposable (elle peut être produite et discutée en justice).
Ces trois exigences se tendent~: plus on protège une pièce, moins elle est
facile à produire et à discuter. Ce cadre, formalisé par l'auteur et ses
collègues~\cite{minka2025cro}, sert de fil rouge aux UE d'investigation
numérique SEC4031 et SEC4052.
\end{boxinfo}

% ============================================================================
\section{L'Entreprise, l'Enquêteur, l'Expert}
\label{sec:c24-roles}
% ============================================================================

\begin{boxjuris}
\textbf{\loicam{n°~2010/012}} (texte officiel)

\medskip
\textbf{Article~52, alinéa~3}~: \textit{« Les Officiers de Police
Judiciaire et les agents habilités de l'Agence peuvent, lors des
investigations, accéder aux moyens de transport, à tout local à usage
professionnel, à l'exclusion des domiciles privés, en vue de rechercher,
de constater les infractions, de demander la communication de tous les
documents professionnels et en prendre copie, recueillir, sur convocation
ou sur place, les renseignements et justifications. »}

\medskip
\textbf{Article~53, alinéa~1}~: \textit{« Les perquisitions en matière de
cybercriminalité sont susceptibles de porter sur des données qui peuvent
être des supports physiques ou des copies réalisées en présence des
personnes qui assistent à la perquisition. »}

\medskip
\textbf{Analyse.} L'enquête appartient aux officiers de police judiciaire
et aux agents habilités. L'entreprise victime préserve et remet~; elle
n'enquête pas sur des personnes et ne qualifie pas les faits
(Chapitre~\ref{chap:droit}).
\end{boxjuris}

\begin{boxmetier}{Analyste SOC}
\textbf{Préserver d'abord.} Avant toute éradication, l'analyste capture la
mémoire des machines clés, copie les journaux hors des machines et note
les empreintes~; il tient un journal de ses actions, heure par heure.
\end{boxmetier}

\begin{boxmetier}{Administrateur systèmes et réseaux}
\textbf{Ne pas éteindre.} L'administrateur isole une machine compromise
du réseau sans l'éteindre, et attend l'équipe de réponse pour la suite.
\end{boxmetier}

\separateur

\begin{boxresume}
\textbf{Ce qu'il faut retenir de ce chapitre~:}
\begin{itemize}
    \item \textbf{Préserver avant d'éradiquer}~: images, mémoire, journaux,
          empreintes.
    \item \textbf{L'ordre de volatilité} (RFC~3227) commence par la mémoire.
    \item \textbf{L'ISO/IEC~27037}~: identification, collecte, acquisition,
          préservation.
    \item \textbf{La chaîne de conservation} trace chaque détention~; une
          rupture fragilise la preuve.
    \item \textbf{L'enquête} appartient aux officiers de police judiciaire
          (loi n°~2010/012, art.~52 et~53)~; l'entreprise préserve et
          remet.
\end{itemize}
\end{boxresume}

\begin{boxinfo}
\textbf{Points de contrôle de l'élève.} Je sais~: expliquer la tension
entre réponse et preuve~; citer l'ordre de volatilité~; nommer les étapes
de l'ISO/IEC~27037~; acquérir un disque avec un bloqueur et des
empreintes~; tenir une chaîne de conservation~; dire qui enquête.
\end{boxinfo}

\bigskip

\begin{boxexo}[title={\ding{45}\quad Exercice~24.1. Volatilité \niveaubase}]
Classez de la plus à la moins volatile~: un journal envoyé au SIEM~; la
liste des connexions réseau~; un fichier sur disque~; le contenu de la
mémoire.
\end{boxexo}

\begin{boxexo}[title={\ding{45}\quad Exercice~24.2. Les étapes \niveaubase}]
Rangez dans les étapes de l'ISO/IEC~27037~: calculer l'empreinte d'une
image~; repérer qu'un routeur garde des journaux~; mettre un disque sous
scellé.
\end{boxexo}

\begin{boxexo}[title={\ding{45}\quad Exercice~24.3. Éteindre~? \niveaubase}]
Pourquoi isoler une machine du réseau plutôt que l'éteindre~? Dans quel
cas faut-il quand même l'éteindre~?
\end{boxexo}

\begin{boxexo}[title={\ding{45}\quad Exercice~24.4. L'empreinte \niveaumoyen}]
L'empreinte de l'image diffère de celle du disque source. Que faites-vous,
et que consignez-vous~?
\end{boxexo}

\begin{boxexo}[title={\ding{45}\quad Exercice~24.5. La fiche \niveaumoyen}]
Rédigez la fiche de chaîne de conservation d'un serveur de Garoua, de sa
saisie à sa remise aux enquêteurs.
\end{boxexo}

\begin{boxexo}[title={\ding{45}\quad Exercice~24.6. Réponse et preuve \niveaumoyen}]
Le directeur veut réinstaller les serveurs dès le samedi matin. Le RSSI
veut d'abord les images. Proposez un compromis qui serve les deux
objectifs.
\end{boxexo}

\begin{boxexo}[title={\ding{45}\quad Exercice~24.7. Les enquêteurs \niveaumoyen}]
Des enquêteurs demandent « tous les journaux ». Que vérifie LiZbiZ, que
remet-elle, et comment le trace-t-elle~?
\end{boxexo}

\begin{boxexo}[title={\ding{45}\quad Exercice~24.8. CRO \niveaumoyen}]
Donnez un exemple où protéger la confidentialité d'une pièce rend plus
difficile son opposabilité.
\end{boxexo}

\begin{boxexo}[title={\ding{45}\quad Exercice~24.9. Le kit \niveauexpert}]
Composez le kit de première réponse de LiZbiZ~: matériel, logiciels,
formulaires, contacts, et la procédure qui l'accompagne.
\end{boxexo}

\begin{boxexo}[title={\ding{45}\quad Exercice~24.10. Réécrire la nuit \niveauexpert}]
Réécrivez la nuit du 5 au 6~mars avec les gestes de ce chapitre et du
Chapitre~\ref{chap:operations}, heure par heure, à partir de 22~h~05.
\end{boxexo}

\bigskip

\begin{boxinfo}
\textbf{Pour aller plus loin.}
\begin{itemize}
    \item ISO/IEC~27037~\cite{iso27037}~; RFC~3227~\cite{rfc3227}~; NIST
          SP~800-86~\cite{nist80086}.
    \item Les UE SEC4031 (Introduction à l'investigation numérique) et
          SEC4052 (Investigation numérique avancée).
\end{itemize}
\end{boxinfo}

\begin{boxinfo}
\textbf{Transition.} Les attaques et les défenses évoluent avec les
technologies. Le Chapitre~\ref{chap:emergents} aborde l'intelligence
artificielle, les objets connectés et la cryptographie post-quantique.
\end{boxinfo}
